\documentclass[12pt,a4paper]{article}
\usepackage[utf8]{inputenc}
\usepackage[T1]{fontenc}
\usepackage{lmodern}
\usepackage{geometry}
\geometry{a4paper, margin=1in}
\usepackage{graphicx}
\usepackage{hyperref}
\usepackage{booktabs}
\usepackage{listings}
\usepackage{xcolor}
\usepackage{float}

\title{
    \vspace{2cm}
    \textbf{\huge University Multimodal RAG Chatbot} \\
    \vspace{0.5cm}
    \large \textit{Project Specifications \& Technical Documentation}
}
\author{Lead AI Engineer}
\date{\today}

\begin{document}

\maketitle
\newpage
\tableofcontents
\newpage

\section{Executive Summary}
This document outlines the detailed architecture, pipeline mechanisms, and deployment instructions for the University Multimodal Retrieval-Augmented Generation (RAG) Chatbot. The system integrates advanced Large Language Models (LLMs), dense vector quantization, real-time embeddings, and multimodal parsing pipelines to serve diverse knowledge-based inquiries.

\section{System Architecture}
The system revolves around a loosely coupled microservice architecture containerized via Docker.
\begin{itemize}
    \item \textbf{API Gateway:} FastAPI asynchronous REST server handling stateless requests.
    \item \textbf{Vector Store:} Qdrant vector database managing separated semantic and visual spaces.
    \item \textbf{Memory Store:} Redis cache handling conversation histories (up to $k=10$ turns).
    \item \textbf{Model Engine:} Abstracted LLM factory utilizing local Ollama deployments or fallback cloud endpoints.
\end{itemize}

\section{Data Pipeline \& Ingestion}
The \texttt{data\_pipeline} module supports hybrid format parsing:
\begin{itemize}
    \item \textbf{Text Formats:} PDF (via PyMuPDF), Word Documents, and textual content dynamically fetched and cleaned from URLs.
    \item \textbf{Images:} Visual data pushed through CLIP (\textit{ViT-B-32}) models to produce 512-dimension image vectors, aligned with textual concepts.
    \item \textbf{Delta Sync Strategy:} Content is hashed (SHA-256). Newly observed documents bypass redundant computation; updates are injected directly into Qdrant nodes preventing the need for LLM fine-tuning.
    \item \textbf{Chunker:} Implements semantic boundary strategies fallback to overlapping logical splitters (size 512, overlap 64) for precise retrieval units.
\end{itemize}

\section{RAG Pipeline Mechanisms}
Retrieval is fully abstract and contextually aware:
\begin{enumerate}
    \item \textbf{Dense Fetching:} Produces similarity candidates using \texttt{nomic-embed-text} for textual queries, and \texttt{CLIP} for image inputs.
    \item \textbf{Reranking Phase:} Initial retrieved limits are re-scored against the specific prompt using a Cross-Encoder (\texttt{ms-marco-MiniLM-L-6-v2}) validating lexical and semantic correlations.
    \item \textbf{Augmented Generation:} Processed context injections limit LLM boundaries to factual constraints, enforcing verifiable source citations.
\end{enumerate}

\section{Benchmark Capabilities}
A built-in benchmark harness utilizes logic inspired by RAGAS. The framework evaluates simulated conversational traces across:
\begin{itemize}
    \item Context Precision \& Context Recall
    \item Answer Relevancy \& Faithfulness
\end{itemize}
Combined weighted output informs developers regarding alignment drift among local models (\texttt{Mistral}, \texttt{Llama 3.1}) or cloud variants (\texttt{GPT-4o}).

\section{Deployment \& Supervisor Management}
To ensure robust scalability on supervisor networks, deployment occurs strictly through exported Docker imagery. 
\begin{itemize}
    \item Build processes create isolated tarballs representing the application state via automated deployment scripts.
    \item Supervisor nodes leverage loading scripts to hydrate local Docker environments detached from public network delays.
    \item Environments can be seamlessly synchronized across production limits enforcing robustness safely offline.
\end{itemize}

\end{document}
